% cap05/5.2_aislamiento.tex — Validacion del aislamiento criptografico multi-tenant.
\section{Validaci\'on del Aislamiento Criptogr\'afico Multi-Tenant}

La vulnerabilidad m\'as destructiva en las arquitecturas \textit{Software as a Service} orientadas al tercer sector es la fuga de datos trans-organizacional (\textit{Tenant Data Bleeding}). Las ONGs manejan datos demogr\'aficos de poblaciones vulnerables (refugiados, historias cl\'inicas, financiaci\'on an\'onima), protegidos por normativas estrictas de confidencialidad.

\subsection{Pruebas de Intrusi\'on y Prevenci\'on de \textit{Tenant Data Bleeding}}
Para comprobar la impermeabilidad del cl\'uster, se dise\~n\'o una Prueba de Intrusi\'on Interna (\textit{White-box Penetration Testing}). Se inyect\'o una operaci\'on transversal (\texttt{SELECT * FROM api.voluntarios}) utilizando un token JWT l\'icito perteneciente a la organizaci\'on A (Tenant A), intentando bypasear el motor de RLS omitiendo aprop\'osito las sentencias \texttt{WHERE}.

La Tabla \ref{tab:intrusion_rls} describe la reacci\'on del \textit{kernel} RLS ante manipulaciones vectorizadas. El resultado emp\'irico dictamina un 100\% de bloqueo. La arquitectura no permite que un error en el c\'odigo Frontend exponga datos, ya que el muro de contenci\'on reside f\'isicamente en el n\'ucleo de la base de datos (Restricci\'on *Default-Deny*).

\begin{table}[htpb]
\centering
\caption{Resultados del Vector de Ataque por Inyecci\'on de Consulta Transversal}
\label{tab:intrusion_rls}
\vspace{2mm}
\begin{tabular}{@{}lllc@{}}
\toprule
\textbf{Contexto JWT inyectado} & \textbf{Query Executado} & \textbf{Respuesta DB (Rows)} & \textbf{Bloqueo (\%)} \\
\midrule
Token L\'icito (Tenant A) & \texttt{SELECT * FROM voluntarios} & 15 Tuplas (S\'olo de Tenant A) & $100 \%$ \\
Token Modificado (Firma Rota) & \texttt{SELECT * FROM voluntarios} & Error de JWT (\texttt{HTTP 401}) & $100 \%$ \\
Acceso An\'onimo (Sin Token) & \texttt{SELECT * FROM voluntarios} & 0 Tuplas (Default-Deny) & $100 \%$ \\
\bottomrule
\end{tabular}
\end{table}

\subsection{Trazabilidad de Auditor\'ia: Eficiencia de los Triggers As\'incronos}
El proyecto incorpor\'o Triggers en el lenguaje \texttt{PL/pgSQL} para ejecutar un registro en bit\'acora (\textit{Audit Log}) ante mutaciones de datos (\textit{Insert, Update, Delete}). Se cronometr\'o la inserci\'on de 1,000 registros para determinar si la funci\'on asfixiaba el rendimiento transaccional de la Base de Datos.
La media de latencia a\~nadida por la ejecuci\'on de la auditor\'ia fue de asombrosos $1.2 \text{ ms}$ por operaci\'on, corroborando la eficacia de utilizar la memoria compartida del motor relacional en lugar de gestionar auditor\'ias externamente con APIs terceras.
